\section{Conclusion and Future Work}
\label{sec:conclusion}

We investigated whether a lightweight linear probe over a small,
decoupled model's hidden states -- extracted without sampling or
generating any text -- can detect specification ambiguity in coding
prompts and rank candidate clarifying questions for an agentic
clarify-then-code pipeline, as a lower-cost alternative to uncertainty
estimation via repeated sampling of the coding model itself. Evaluated
against a zero-shot self-rank baseline under grouped, leave-one-source-out
splits across four benchmark sources, the probe reaches ranking quality
close to, and on its training-distribution source above, the baseline, at
substantially lower measured per-decision inference cost, framing this as
a cost-quality tradeoff rather than an unconditional accuracy improvement.
Deployed in a live, two-round elicit-rank-clarify-code loop against real
executable test suites, we found that ranking the labeled ambiguity
correctly is associated with a statistically significant round-1
correctness gain when measured by fractional test-case pass rate, though
not when measured by an all-or-nothing pass metric on the same data -- a
discrepancy we treat as evidence that correctness-metric choice materially
affects what an agentic clarification pipeline's evaluation can detect.
We further found that asking any clarifying question at all, independent
of ranking precision, is itself a robust and significant correctness
improvement, and that a cheap, targeted prompt augmentation addressing a
separate, orthogonal source of crash-class failures further improves
outcomes.

We reported several points where earlier-stage conclusions in this
investigation -- most notably an initial finding that the probe
outperformed zero-shot ranking -- did not survive evaluation at larger
scale and under stricter filtering, and disclose these reversals as part
of the record rather than reporting only the final numbers. We believe
this level of disclosure is necessary for agentic software engineering
research broadly, where pipelines with several interacting stages (elicit,
rank, clarify, generate, test) create many opportunities for an early,
underpowered, or confounded result to look like a finding.

Future work includes: benchmarking directly against a repeated-sampling,
self-consistency baseline on identical problems to substantiate the
cost-efficiency argument empirically rather than by citation; measuring
the actual batched $N$-candidate extraction cost against zero-shot
generation, rather than the bare-prompt-only forward pass reported here;
characterizing the tradeoff between probe-host model size and both ranking
quality and inference cost; and improving cross-source generalization for
long, free-form prompts where the ambiguous edit constitutes a small
fraction of the total prompt, which we identified as the ranking probe's
principal weakness.
